\chapter{{\fontsize{14pt}{21pt}\selectfont\MakeUppercase{Предметная область}}}
\label{ch:subject-area}

\section{Анализ предметной области (исследовательская часть)}

Цифровое присутствие локального бизнеса определяется как совокупность доступных в~цифровых каналах данных и сигналов, по которым потенциальный клиент формирует ожидания о~компании и принимает решение о~взаимодействии. Ядро составляют <<карточки организации>> (профили) на платформах бизнес-информации~--- сервисах, где отображаются адрес, часы работы, контакты и другие атрибуты компании~\cite{src:4, src:3}. Помимо профилей на картографических сервисах, компоненты цифрового присутствия включают страницы и сообщества в~социальных сетях (VK), каналы и боты в~мессенджерах (Telegram), агрегаторы товаров и услуг, отзывы и рейтинги на различных платформах.

Множественность каналов формирует задачу поддержания единообразной и актуальной информации о~компании на каждом из них~--- задачу обеспечения \textit{консистентности данных} цифрового присутствия. Для потребителей критичны базовые сведения: часы работы и контактно-адресная информация. В~Российской Федерации пользователи геосервисов часто используют онлайн-карты для уточнения графика работы (48~\%)~\cite{src:35}; при выборе мест обращают внимание на рейтинг (59~\%) и отзывы (51~\%)~\cite{src:35}. Таким образом, исследовательская часть проекта охватывает:

\begin{enumerate}
  \item консистентность справочных данных между платформами;
  \item управление отзывами как репутационным сигналом;
  \item классификацию целевых платформ по типу интеграции (API / RPA).
\end{enumerate}

В~российском контексте три геосервиса делят рынок: Яндекс Карты~--- 53~\%, 2ГИС~--- 44~\%, Google Карты~--- 3~\%~\cite{src:35}. При этом ни Яндекс.Бизнес, ни 2ГИС не предоставляют публичного API для управления карточкой организации, что требует гибридного подхода к~проектированию интеграций. Классификация платформ по типу интеграции представлена в~таблице~\ref{tab:platform-class}.

\begin{table}[H]
  \caption{Классификация целевых платформ по типу интеграции}
  \label{tab:platform-class}
  \small
  \begin{tabularx}{\textwidth}{|p{3.5cm}|p{2.5cm}|p{2cm}|X|}
    \hline
    \textbf{Платформа} & \textbf{Метод} & \textbf{Приоритет} & \textbf{Примечание} \\
    \hline
    VK & API (официальный) & MVP & OAuth~2.0, методы для сообществ \\
    \hline
    Telegram & API (Bot API) & MVP & Токен бота, HTTP-API \\
    \hline
    Яндекс.Бизнес & RPA (Playwright) & MVP & Нет публичного API, браузерная автоматизация \\
    \hline
    2ГИС & RPA (Playwright) & Фаза~2 & Нет публичного API \\
    \hline
    Google Business Profile & API (официальный) & Фаза~3 & 3~\% использования в~РФ \\
    \hline
  \end{tabularx}
\end{table}
